% =============================================================================
% Finance Academic Paper — Template
% =============================================================================
% Usage: copy into drafts/documents/ and rename, then edit.
% Compile: cd drafts/documents
%   TEXINPUTS=../latex_files:$TEXINPUTS pdflatex -interaction=nonstopmode file.tex
%   BIBINPUTS=../latex_files:$BIBINPUTS BSTINPUTS=../latex_files:$BSTINPUTS bibtex file
%   (then two more pdflatex passes)
% =============================================================================

\documentclass[letterpaper,11pt]{article}
% =============================================================================
% Pathways: Economics & Finance Workshop — Document Preamble
% =============================================================================
% Usage: Each handout/activity .tex file should start with:
%   \documentclass[letterpaper,10pt]{article}
%   % =============================================================================
% Pathways: Economics & Finance Workshop — Document Preamble
% =============================================================================
% Usage: Each handout/activity .tex file should start with:
%   \documentclass[letterpaper,10pt]{article}
%   % =============================================================================
% Pathways: Economics & Finance Workshop — Document Preamble
% =============================================================================
% Usage: Each handout/activity .tex file should start with:
%   \documentclass[letterpaper,10pt]{article}
%   \input{../../codes/Preambles/header_doc}
% Compile: cd output/Handouts && TEXINPUTS=../../codes/Preambles:$TEXINPUTS xelatex file.tex
% =============================================================================

% ---------- Packages ----------
% Fonts
\usepackage{palatino}
\usepackage[utf8]{inputenc}
\usepackage[T1]{fontenc}
\usepackage[normalem]{ulem}

% Math
\usepackage{amsmath,amsthm}
\usepackage{newpxmath}
\usepackage{bm}
\usepackage{mathrsfs}
\usepackage{dsfont}

% Figures and tables
\usepackage{graphicx}
\usepackage{placeins}
\usepackage{xcolor}
\usepackage{array,tabularx,booktabs,longtable,subcaption}
\usepackage{threeparttable}
\usepackage{pdflscape}
\usepackage{siunitx}

% TikZ and pgfplots
\usepackage{pgfplots}
\pgfplotsset{compat=1.18}
\usepackage{tikz}
\usetikzlibrary{positioning,arrows.meta,calc}

% Colored boxes
\usepackage[most]{tcolorbox}

% Citations (natbib loaded before hyperref to avoid conflicts)
\usepackage{natbib}

% Miscellaneous
\usepackage{hyperref}
% Links stay clickable but render as plain text: `hidelinks` removes BOTH the
% coloured border boxes (hyperref's default for links/citations/URLs) and any
% link colouring. Must come after the package load.
\hypersetup{hidelinks}
\usepackage{setspace}
\usepackage{enumitem}

% ---------- Margins ----------
\usepackage[left=1in,right=1in,top=0.75in,bottom=1in,marginparwidth=1.5in,footskip=0.7in]{geometry}

% ---------- Itemize ----------
\setlist[enumerate]{itemsep=-0.4em,topsep=0.2em}
\setlist[itemize]{itemsep=-0.4em,topsep=0.2em}

% =============================================================================
% Colors — Okabe-Ito colorblind-friendly palette
% =============================================================================

\definecolor{black}{HTML}{000000}
\definecolor{orange}{HTML}{E69F00}
\definecolor{skyblue}{HTML}{56B4E9}
\definecolor{green}{HTML}{009E73}
\definecolor{yellow}{HTML}{F0E442}
\definecolor{blue}{HTML}{0072B2}
\definecolor{vermillion}{HTML}{D55E00}
\definecolor{purple}{HTML}{CC79A7}

% Semantic aliases (mapped to Okabe-Ito)
\definecolor{softbg}{HTML}{F5F5F5}
\definecolor{lightgray}{gray}{0.65}

% =============================================================================
% Box Environments
% =============================================================================

% --- Key points (orange background) ---
\newtcolorbox{keybox}{
  colback=orange!12,
  colframe=orange,
  fonttitle=\bfseries,
  boxrule=0.5pt,
  arc=2pt,
  left=6pt,
  right=6pt,
  top=4pt,
  bottom=4pt
}

% --- Highlights (orange left accent) ---
\newtcolorbox{highlightbox}{
  colback=white,
  colframe=orange,
  fonttitle=\bfseries,
  boxrule=0pt,
  leftrule=3pt,
  arc=0pt,
  left=6pt,
  right=6pt,
  top=4pt,
  bottom=4pt
}

% --- Definitions (blue border, titled) ---
\newtcolorbox{definitionbox}[1][]{
  colback=blue!5,
  colframe=blue,
  fonttitle=\bfseries,
  title={#1},
  boxrule=0.5pt,
  arc=2pt,
  left=6pt,
  right=6pt,
  top=4pt,
  bottom=4pt
}

% --- Methods (sky blue background) ---
\newtcolorbox{methodbox}{
  colback=skyblue!8,
  colframe=skyblue!70,
  fonttitle=\bfseries,
  boxrule=0.5pt,
  arc=2pt,
  left=6pt,
  right=6pt,
  top=4pt,
  bottom=4pt
}

% --- Results (green border) ---
\newtcolorbox{resultbox}{
  colback=green!5,
  colframe=green,
  fonttitle=\bfseries,
  boxrule=0.5pt,
  arc=2pt,
  left=6pt,
  right=6pt,
  top=4pt,
  bottom=4pt
}

% --- Quotes (purple left rule) ---
\newtcolorbox{quotebox}{
  colback=softbg,
  colframe=purple,
  fonttitle=\itshape,
  boxrule=0pt,
  leftrule=2pt,
  arc=0pt,
  left=6pt,
  right=6pt,
  top=4pt,
  bottom=4pt
}

% =============================================================================
% Math Commands — Notation Legend
% =============================================================================

% --- Spaces ---
\newcommand{\R}{\mathbb{R}}
\newcommand{\Rplus}{\mathbb{R}_+}
\newcommand{\Rk}{\mathbb{R}^k}

% --- Probability ---
\newcommand{\Prob}{\mathbb{P}}

% --- Expectation and moments ---
\newcommand{\E}{\mathbb{E}}
\DeclareMathOperator{\var}{var}
\DeclareMathOperator{\cov}{cov}
\DeclareMathOperator{\corr}{corr}
\newcommand{\BLP}{\mathscr{P}}

% --- Convergence ---
\DeclareMathOperator{\plim}{plim}
\newcommand{\convp}{\xrightarrow{\quad p \quad}}
\newcommand{\convd}{\xrightarrow{\quad d \quad}}

% --- Likelihood and information ---
\newcommand{\Lcal}{\mathcal{L}}
\newcommand{\Infmat}{\mathscr{I}}

% --- Relations and operators ---
\newcommand{\defeq}{\overset{\text{def}}{=}}
\newcommand{\ind}{\mathds{1}}
\newcommand{\Normal}{\mathrm{N}}

% --- Calculus shortcuts ---
\newcommand{\suminfty}[1]{\sum_{#1=0}^\infty}
\newcommand{\prtl}[2]{\frac{\partial #1}{\partial #2}}
\renewcommand{\div}[2]{\frac{d #1}{d #2}}
\newcommand{\suchthat}{\text{ s.t. }}
\newcommand{\wrt}{\text{ w.r.t. }}

% --- Text color shortcuts ---
\newcommand{\red}[1]{\textcolor{vermillion}{#1}}
\newcommand{\blue}[1]{\textcolor{blue}{#1}}
\newcommand{\pmark}{\textcolor{green}{\checkmark}}
\newcommand{\xmark}{\textcolor{vermillion}{\texttimes}}

% --- Stata esttab significance stars ---
\newcommand{\sym}[1]{\ensuremath{^{#1}}}

% Compile: cd output/Handouts && TEXINPUTS=../../codes/Preambles:$TEXINPUTS xelatex file.tex
% =============================================================================

% ---------- Packages ----------
% Fonts
\usepackage{palatino}
\usepackage[utf8]{inputenc}
\usepackage[T1]{fontenc}
\usepackage[normalem]{ulem}

% Math
\usepackage{amsmath,amsthm}
\usepackage{newpxmath}
\usepackage{bm}
\usepackage{mathrsfs}
\usepackage{dsfont}

% Figures and tables
\usepackage{graphicx}
\usepackage{placeins}
\usepackage{xcolor}
\usepackage{array,tabularx,booktabs,longtable,subcaption}
\usepackage{threeparttable}
\usepackage{pdflscape}
\usepackage{siunitx}

% TikZ and pgfplots
\usepackage{pgfplots}
\pgfplotsset{compat=1.18}
\usepackage{tikz}
\usetikzlibrary{positioning,arrows.meta,calc}

% Colored boxes
\usepackage[most]{tcolorbox}

% Citations (natbib loaded before hyperref to avoid conflicts)
\usepackage{natbib}

% Miscellaneous
\usepackage{hyperref}
% Links stay clickable but render as plain text: `hidelinks` removes BOTH the
% coloured border boxes (hyperref's default for links/citations/URLs) and any
% link colouring. Must come after the package load.
\hypersetup{hidelinks}
\usepackage{setspace}
\usepackage{enumitem}

% ---------- Margins ----------
\usepackage[left=1in,right=1in,top=0.75in,bottom=1in,marginparwidth=1.5in,footskip=0.7in]{geometry}

% ---------- Itemize ----------
\setlist[enumerate]{itemsep=-0.4em,topsep=0.2em}
\setlist[itemize]{itemsep=-0.4em,topsep=0.2em}

% =============================================================================
% Colors — Okabe-Ito colorblind-friendly palette
% =============================================================================

\definecolor{black}{HTML}{000000}
\definecolor{orange}{HTML}{E69F00}
\definecolor{skyblue}{HTML}{56B4E9}
\definecolor{green}{HTML}{009E73}
\definecolor{yellow}{HTML}{F0E442}
\definecolor{blue}{HTML}{0072B2}
\definecolor{vermillion}{HTML}{D55E00}
\definecolor{purple}{HTML}{CC79A7}

% Semantic aliases (mapped to Okabe-Ito)
\definecolor{softbg}{HTML}{F5F5F5}
\definecolor{lightgray}{gray}{0.65}

% =============================================================================
% Box Environments
% =============================================================================

% --- Key points (orange background) ---
\newtcolorbox{keybox}{
  colback=orange!12,
  colframe=orange,
  fonttitle=\bfseries,
  boxrule=0.5pt,
  arc=2pt,
  left=6pt,
  right=6pt,
  top=4pt,
  bottom=4pt
}

% --- Highlights (orange left accent) ---
\newtcolorbox{highlightbox}{
  colback=white,
  colframe=orange,
  fonttitle=\bfseries,
  boxrule=0pt,
  leftrule=3pt,
  arc=0pt,
  left=6pt,
  right=6pt,
  top=4pt,
  bottom=4pt
}

% --- Definitions (blue border, titled) ---
\newtcolorbox{definitionbox}[1][]{
  colback=blue!5,
  colframe=blue,
  fonttitle=\bfseries,
  title={#1},
  boxrule=0.5pt,
  arc=2pt,
  left=6pt,
  right=6pt,
  top=4pt,
  bottom=4pt
}

% --- Methods (sky blue background) ---
\newtcolorbox{methodbox}{
  colback=skyblue!8,
  colframe=skyblue!70,
  fonttitle=\bfseries,
  boxrule=0.5pt,
  arc=2pt,
  left=6pt,
  right=6pt,
  top=4pt,
  bottom=4pt
}

% --- Results (green border) ---
\newtcolorbox{resultbox}{
  colback=green!5,
  colframe=green,
  fonttitle=\bfseries,
  boxrule=0.5pt,
  arc=2pt,
  left=6pt,
  right=6pt,
  top=4pt,
  bottom=4pt
}

% --- Quotes (purple left rule) ---
\newtcolorbox{quotebox}{
  colback=softbg,
  colframe=purple,
  fonttitle=\itshape,
  boxrule=0pt,
  leftrule=2pt,
  arc=0pt,
  left=6pt,
  right=6pt,
  top=4pt,
  bottom=4pt
}

% =============================================================================
% Math Commands — Notation Legend
% =============================================================================

% --- Spaces ---
\newcommand{\R}{\mathbb{R}}
\newcommand{\Rplus}{\mathbb{R}_+}
\newcommand{\Rk}{\mathbb{R}^k}

% --- Probability ---
\newcommand{\Prob}{\mathbb{P}}

% --- Expectation and moments ---
\newcommand{\E}{\mathbb{E}}
\DeclareMathOperator{\var}{var}
\DeclareMathOperator{\cov}{cov}
\DeclareMathOperator{\corr}{corr}
\newcommand{\BLP}{\mathscr{P}}

% --- Convergence ---
\DeclareMathOperator{\plim}{plim}
\newcommand{\convp}{\xrightarrow{\quad p \quad}}
\newcommand{\convd}{\xrightarrow{\quad d \quad}}

% --- Likelihood and information ---
\newcommand{\Lcal}{\mathcal{L}}
\newcommand{\Infmat}{\mathscr{I}}

% --- Relations and operators ---
\newcommand{\defeq}{\overset{\text{def}}{=}}
\newcommand{\ind}{\mathds{1}}
\newcommand{\Normal}{\mathrm{N}}

% --- Calculus shortcuts ---
\newcommand{\suminfty}[1]{\sum_{#1=0}^\infty}
\newcommand{\prtl}[2]{\frac{\partial #1}{\partial #2}}
\renewcommand{\div}[2]{\frac{d #1}{d #2}}
\newcommand{\suchthat}{\text{ s.t. }}
\newcommand{\wrt}{\text{ w.r.t. }}

% --- Text color shortcuts ---
\newcommand{\red}[1]{\textcolor{vermillion}{#1}}
\newcommand{\blue}[1]{\textcolor{blue}{#1}}
\newcommand{\pmark}{\textcolor{green}{\checkmark}}
\newcommand{\xmark}{\textcolor{vermillion}{\texttimes}}

% --- Stata esttab significance stars ---
\newcommand{\sym}[1]{\ensuremath{^{#1}}}

% Compile: cd output/Handouts && TEXINPUTS=../../codes/Preambles:$TEXINPUTS xelatex file.tex
% =============================================================================

% ---------- Packages ----------
% Fonts
\usepackage{palatino}
\usepackage[utf8]{inputenc}
\usepackage[T1]{fontenc}
\usepackage[normalem]{ulem}

% Math
\usepackage{amsmath,amsthm}
\usepackage{newpxmath}
\usepackage{bm}
\usepackage{mathrsfs}
\usepackage{dsfont}

% Figures and tables
\usepackage{graphicx}
\usepackage{placeins}
\usepackage{xcolor}
\usepackage{array,tabularx,booktabs,longtable,subcaption}
\usepackage{threeparttable}
\usepackage{pdflscape}
\usepackage{siunitx}

% TikZ and pgfplots
\usepackage{pgfplots}
\pgfplotsset{compat=1.18}
\usepackage{tikz}
\usetikzlibrary{positioning,arrows.meta,calc}

% Colored boxes
\usepackage[most]{tcolorbox}

% Citations (natbib loaded before hyperref to avoid conflicts)
\usepackage{natbib}

% Miscellaneous
\usepackage{hyperref}
% Links stay clickable but render as plain text: `hidelinks` removes BOTH the
% coloured border boxes (hyperref's default for links/citations/URLs) and any
% link colouring. Must come after the package load.
\hypersetup{hidelinks}
\usepackage{setspace}
\usepackage{enumitem}

% ---------- Margins ----------
\usepackage[left=1in,right=1in,top=0.75in,bottom=1in,marginparwidth=1.5in,footskip=0.7in]{geometry}

% ---------- Itemize ----------
\setlist[enumerate]{itemsep=-0.4em,topsep=0.2em}
\setlist[itemize]{itemsep=-0.4em,topsep=0.2em}

% =============================================================================
% Colors — Okabe-Ito colorblind-friendly palette
% =============================================================================

\definecolor{black}{HTML}{000000}
\definecolor{orange}{HTML}{E69F00}
\definecolor{skyblue}{HTML}{56B4E9}
\definecolor{green}{HTML}{009E73}
\definecolor{yellow}{HTML}{F0E442}
\definecolor{blue}{HTML}{0072B2}
\definecolor{vermillion}{HTML}{D55E00}
\definecolor{purple}{HTML}{CC79A7}

% Semantic aliases (mapped to Okabe-Ito)
\definecolor{softbg}{HTML}{F5F5F5}
\definecolor{lightgray}{gray}{0.65}

% =============================================================================
% Box Environments
% =============================================================================

% --- Key points (orange background) ---
\newtcolorbox{keybox}{
  colback=orange!12,
  colframe=orange,
  fonttitle=\bfseries,
  boxrule=0.5pt,
  arc=2pt,
  left=6pt,
  right=6pt,
  top=4pt,
  bottom=4pt
}

% --- Highlights (orange left accent) ---
\newtcolorbox{highlightbox}{
  colback=white,
  colframe=orange,
  fonttitle=\bfseries,
  boxrule=0pt,
  leftrule=3pt,
  arc=0pt,
  left=6pt,
  right=6pt,
  top=4pt,
  bottom=4pt
}

% --- Definitions (blue border, titled) ---
\newtcolorbox{definitionbox}[1][]{
  colback=blue!5,
  colframe=blue,
  fonttitle=\bfseries,
  title={#1},
  boxrule=0.5pt,
  arc=2pt,
  left=6pt,
  right=6pt,
  top=4pt,
  bottom=4pt
}

% --- Methods (sky blue background) ---
\newtcolorbox{methodbox}{
  colback=skyblue!8,
  colframe=skyblue!70,
  fonttitle=\bfseries,
  boxrule=0.5pt,
  arc=2pt,
  left=6pt,
  right=6pt,
  top=4pt,
  bottom=4pt
}

% --- Results (green border) ---
\newtcolorbox{resultbox}{
  colback=green!5,
  colframe=green,
  fonttitle=\bfseries,
  boxrule=0.5pt,
  arc=2pt,
  left=6pt,
  right=6pt,
  top=4pt,
  bottom=4pt
}

% --- Quotes (purple left rule) ---
\newtcolorbox{quotebox}{
  colback=softbg,
  colframe=purple,
  fonttitle=\itshape,
  boxrule=0pt,
  leftrule=2pt,
  arc=0pt,
  left=6pt,
  right=6pt,
  top=4pt,
  bottom=4pt
}

% =============================================================================
% Math Commands — Notation Legend
% =============================================================================

% --- Spaces ---
\newcommand{\R}{\mathbb{R}}
\newcommand{\Rplus}{\mathbb{R}_+}
\newcommand{\Rk}{\mathbb{R}^k}

% --- Probability ---
\newcommand{\Prob}{\mathbb{P}}

% --- Expectation and moments ---
\newcommand{\E}{\mathbb{E}}
\DeclareMathOperator{\var}{var}
\DeclareMathOperator{\cov}{cov}
\DeclareMathOperator{\corr}{corr}
\newcommand{\BLP}{\mathscr{P}}

% --- Convergence ---
\DeclareMathOperator{\plim}{plim}
\newcommand{\convp}{\xrightarrow{\quad p \quad}}
\newcommand{\convd}{\xrightarrow{\quad d \quad}}

% --- Likelihood and information ---
\newcommand{\Lcal}{\mathcal{L}}
\newcommand{\Infmat}{\mathscr{I}}

% --- Relations and operators ---
\newcommand{\defeq}{\overset{\text{def}}{=}}
\newcommand{\ind}{\mathds{1}}
\newcommand{\Normal}{\mathrm{N}}

% --- Calculus shortcuts ---
\newcommand{\suminfty}[1]{\sum_{#1=0}^\infty}
\newcommand{\prtl}[2]{\frac{\partial #1}{\partial #2}}
\renewcommand{\div}[2]{\frac{d #1}{d #2}}
\newcommand{\suchthat}{\text{ s.t. }}
\newcommand{\wrt}{\text{ w.r.t. }}

% --- Text color shortcuts ---
\newcommand{\red}[1]{\textcolor{vermillion}{#1}}
\newcommand{\blue}[1]{\textcolor{blue}{#1}}
\newcommand{\pmark}{\textcolor{green}{\checkmark}}
\newcommand{\xmark}{\textcolor{vermillion}{\texttimes}}

% --- Stata esttab significance stars ---
\newcommand{\sym}[1]{\ensuremath{^{#1}}}


% ---------- Template-local setup ----------
\usepackage[round]{natbib}          % \citet / \citep for author-year citations
\bibliographystyle{../latex_files/aea}
\def\sym#1{\ifmmode^{#1}\else\(^{#1}\)\fi}   % significance stars in tables

\onehalfspacing

\title{\textbf{Paper Title:\\ A Subtitle That States the Contribution}}
\author{
  Your Name\thanks{%
    Institution, address, email. Acknowledgments, disclosures, and funding go here.
    Any remaining errors are our own.}
  \and
  Coauthor Name\thanks{Institution, email.}
}
\date{\today}

\begin{document}

\maketitle

% =============================================================================
% Abstract
% =============================================================================
\begin{abstract}
\noindent
One sentence on the question and why it matters. One sentence on the data
(source, unit of observation, sample period). One or two sentences on the
identification strategy or empirical design. Two sentences on the main result,
with the magnitude stated in economically interpretable units. One sentence on
the mechanism or interpretation. One sentence on the contribution or
implication for practice and policy.

\vspace{1em}
\noindent \textbf{Keywords:} keyword one, keyword two, keyword three

\noindent \textbf{JEL Classification:} G12, G14, G32
\end{abstract}

\thispagestyle{empty}
\newpage
\setcounter{page}{1}

% =============================================================================
\section{Introduction}
\label{sec:intro}
% =============================================================================

\textit{Hook.} Open with the economic phenomenon and the puzzle it creates.
State why a reader who cares about asset pricing, corporate finance, or
financial intermediation should care about the answer.

\textit{Research question.} State the question precisely, in one sentence, in
terms of the object you actually estimate.

\textit{Setting and data.} Describe the empirical setting and why it is
well-suited to the question. Name the data sources and the sample period.

\textit{Empirical strategy.} Summarize the identification argument in
non-technical terms: what variation identifies the parameter of interest, and
what would have to be true for the estimate to be causal.

\textit{Results.} Preview the main findings with magnitudes. Then preview the
robustness checks and the mechanism evidence that rule out leading alternative
explanations.

\textit{Related literature.} This paper contributes to three strands of the
literature. First, it builds on work on [strand one], which establishes
[finding] \citep{Angrist1996_iv}. That literature takes [assumption] as given;
we relax it by [what you do differently]. Second, it relates to research on
[strand two], where \citet{Callaway2021_did} show [finding]; relative to that
work, our contribution is [the marginal contribution, stated concretely].
Third, it speaks to the applied literature on [strand three], which has
documented [stylized fact] but has not addressed [the gap]. We are, to our
knowledge, the first to [contribution]. \textit{Keep this to a single
paragraph; be explicit about the marginal contribution rather than merely
listing citations.}

\textit{Roadmap.} Section~\ref{sec:data} describes the data.
Section~\ref{sec:method} presents the empirical strategy.
Section~\ref{sec:results} reports the results. Section~\ref{sec:conclusion}
concludes.

% =============================================================================
\section{Data}
\label{sec:data}
% =============================================================================

\subsection{Data Sources}

Describe each source: provider, coverage, frequency, and the identifier used to
merge. For example: security-level returns from CRSP, accounting fundamentals
from Compustat, institutional holdings from 13F filings.

\subsection{Sample Construction}

State the sample period and every filter applied, in the order applied. Report
how many observations each filter removes so the sample is reproducible.

\begin{table}[htbp]
  \centering
  \caption{Sample Construction}
  \label{tab:sample}
  \begin{tabular}{lcc}
    \toprule
    Filter & Observations & Firms \\
    \midrule
    Raw merged panel      & 000{,}000 & 00{,}000 \\
    Drop financials       & 000{,}000 & 00{,}000 \\
    Require non-missing outcome & 000{,}000 & 00{,}000 \\
    \midrule
    Final sample          & 000{,}000 & 00{,}000 \\
    \bottomrule
  \end{tabular}

  \vspace{0.5em}
  {\footnotesize \textit{Notes:} Sample period is YYYY--YYYY. Each row reports
  the count remaining after the stated filter is applied.}
\end{table}

\subsection{Variable Definitions}

Define the outcome, treatment, and controls. Give the exact construction
(including winsorization and any deflators) here or in
Appendix~\ref{app:variables}.

\subsection{Summary Statistics}

\begin{table}[htbp]
  \centering
  \caption{Summary Statistics}
  \label{tab:sumstats}
  \begin{tabular}{lccccc}
    \toprule
    Variable & N & Mean & SD & P25 & P75 \\
    \midrule
    Outcome $y_{it}$    & 000{,}000 & 0.00 & 0.00 & 0.00 & 0.00 \\
    Treatment $D_{it}$  & 000{,}000 & 0.00 & 0.00 & 0.00 & 0.00 \\
    Size (log assets)   & 000{,}000 & 0.00 & 0.00 & 0.00 & 0.00 \\
    Leverage            & 000{,}000 & 0.00 & 0.00 & 0.00 & 0.00 \\
    \bottomrule
  \end{tabular}

  \vspace{0.5em}
  {\footnotesize \textit{Notes:} Continuous variables are winsorized at the 1st
  and 99th percentiles. Variable definitions are in
  Appendix~\ref{app:variables}.}
\end{table}

% =============================================================================
\section{Methodology}
\label{sec:method}
% =============================================================================

\subsection{Empirical Specification}

State the estimating equation and define every term.

\begin{equation}
  y_{it} = \alpha + \beta \, D_{it} + \gamma' X_{it} + \mu_i + \lambda_t + \varepsilon_{it},
  \label{eq:main}
\end{equation}

\noindent where $y_{it}$ is the outcome for firm $i$ in period $t$, $D_{it}$ is
the treatment, $X_{it}$ is a vector of controls, and $\mu_i$ and $\lambda_t$
are firm and time fixed effects. The coefficient of interest is $\beta$.
Standard errors are clustered at the [level]; state why that is the level at
which treatment is assigned.

\subsection{Identification}

\begin{methodbox}
\textbf{Identifying assumption.} State it formally, then in words.
\[
  \E\!\left[\varepsilon_{it} \mid D_{it}, X_{it}, \mu_i, \lambda_t\right] = 0
\]
Explain what would violate it and what evidence you provide against that
violation.
\end{methodbox}

\subsection{Threats to Identification}

List the leading alternative explanations and the test that addresses each:
selection on unobservables, reverse causality, measurement error in the
treatment, and differential pre-trends.

% =============================================================================
\section{Results}
\label{sec:results}
% =============================================================================

\subsection{Main Results}

\begin{table}[htbp]
  \centering
  \caption{Effect of $D$ on $y$}
  \label{tab:main}
  \begin{tabular}{lcccc}
    \toprule
    & (1) & (2) & (3) & (4) \\
    \midrule
    Treatment $D_{it}$ & 0.000\sym{***} & 0.000\sym{***} & 0.000\sym{**} & 0.000\sym{**} \\
                       & (0.000) & (0.000) & (0.000) & (0.000) \\
    \midrule
    Controls           & No  & Yes & Yes & Yes \\
    Firm FE            & No  & No  & Yes & Yes \\
    Year FE            & No  & No  & No  & Yes \\
    Observations       & 000{,}000 & 000{,}000 & 000{,}000 & 000{,}000 \\
    $R^2$              & 0.00 & 0.00 & 0.00 & 0.00 \\
    \bottomrule
  \end{tabular}

  \vspace{0.5em}
  {\footnotesize \textit{Notes:} Standard errors clustered by [level] in
  parentheses. \sym{*} $p<0.10$, \sym{**} $p<0.05$, \sym{***} $p<0.01$.}
\end{table}

Interpret the magnitude in economic terms: a one-standard-deviation increase in
$D$ is associated with a change of X basis points in $y$, or Y\% of the sample
mean. Compare the estimate to benchmarks in the literature.

\subsection{Robustness}

Report alternative specifications, alternative samples, alternative clustering,
and placebo tests. Say what each check rules out.

\subsection{Heterogeneity and Mechanisms}

Split the sample along the dimension implied by the theory. State the sign the
mechanism predicts \textit{before} reporting the split, so the test is
falsifiable.

\begin{figure}[htbp]
  \centering
  % \includegraphics[width=0.8\textwidth]{../../output/figures/event_study.pdf}
  \caption{Event-Study Estimates}
  \label{fig:event}

  \vspace{0.5em}
  {\footnotesize \textit{Notes:} Point estimates and 95\% confidence intervals
  from Equation~\eqref{eq:main} estimated with leads and lags. Period $-1$ is
  the omitted category.}
\end{figure}

% =============================================================================
\section{Conclusion}
\label{sec:conclusion}
% =============================================================================

Restate the question and the answer, with the magnitude. State what the result
implies for investors, firms, or regulators. Be explicit about the limits of
external validity: the setting, the sample period, and the population the
estimate speaks to. Close with the open question your evidence does not settle.

% =============================================================================
% References
% =============================================================================
\newpage
\bibliography{../latex_files/references}

% =============================================================================
% Appendix
% =============================================================================
\newpage
\appendix
\section{Variable Definitions}
\label{app:variables}

Give the exact construction of each variable: source item, formula, units, and
any winsorization or deflation.

\section{Additional Results}
\label{app:additional}

Supporting tables and figures referenced in the main text but not central to
the argument.

\section{Data Appendix}
\label{app:data}

Merge procedure, treatment of missing values, and any deviations from the
standard filters used in the literature.

\end{document}
